\documentclass[aspectratio=169,11pt]{beamer}

\usepackage[T1]{fontenc}
\usepackage{lmodern}
\usepackage{booktabs}

\usetheme{default}
\usecolortheme{default}
\setbeamertemplate{navigation symbols}{}
\setbeamertemplate{caption}[numbered]

% ---------- Color Palette ----------
\definecolor{navy}{RGB}{16,34,66}
\definecolor{gold}{RGB}{176,141,58}
\definecolor{cream}{RGB}{250,248,242}
\definecolor{slate}{RGB}{60,70,85}

\setbeamercolor{normal text}{fg=slate,bg=cream}
\setbeamercolor{structure}{fg=navy}
\setbeamercolor{frametitle}{fg=cream,bg=navy}
\setbeamercolor{title}{fg=cream,bg=navy}
\setbeamercolor{subtitle}{fg=gold}
\setbeamercolor{block title}{fg=cream,bg=navy}
\setbeamercolor{block body}{fg=slate,bg=cream!80!navy!5}
\setbeamercolor{alerted text}{fg=gold!80!black}
\setbeamercolor{itemize item}{fg=gold}
\setbeamercolor{itemize subitem}{fg=navy}

\setbeamerfont{frametitle}{series=\bfseries,size=\large}
\setbeamerfont{title}{series=\bfseries,size=\LARGE}

% ---------- Metadata ----------
\title[Pro Se Mandamus Guide]{How to File a Writ of Mandamus Pro Se}
\subtitle{A Federal Litigation Guide for Immigration Delay Cases}
\author{}
\date{Reviewed July 2026}

\begin{document}
	
	% ---------- Title ----------
	\begin{frame}[plain]
		\titlepage
		\vfill
		\begin{center}
			\small\textcolor{gold}{\textbf{Attorney Advertising} \textperiodcentered\ General legal information, not legal advice}
		\end{center}
	\end{frame}
	
	% ---------- Agenda ----------
	\begin{frame}{Roadmap}
		\begin{columns}[T]
			\begin{column}{0.5\textwidth}
				\begin{enumerate}
					\item Is Mandamus Right for Your Case?
					\item Exhaust Administrative Remedies
					\item Choosing Your District Court
					\item Identifying Defendants
					\item Anatomy of the Complaint
					\item Supporting Documents
				\end{enumerate}
			\end{column}
			\begin{column}{0.5\textwidth}
				\begin{enumerate}\setcounter{enumi}{6}
					\item Filing via PACER / CM-ECF
					\item Filing Fees \& IFP
					\item Serving the United States (Rule 4(i))
					\item The 60-Day Window
					\item Ten Common Mistakes
					\item When to Hire Counsel
				\end{enumerate}
			\end{column}
		\end{columns}
	\end{frame}
	
	% ---------- Read this first ----------
	\begin{frame}{Read This First}
		\begin{block}{A Federal Lawsuit Is Not Simple, But It Is Doable}
			Mandamus is mechanically straightforward: the theory is settled, the facts are short, and the government rarely wants to litigate.
		\end{block}
		\begin{itemize}
			\item FRCP 11(a) and 28 U.S.C. \S~1654 permit individuals to appear pro se.
			\item But federal procedure is unforgiving --- a missed deadline, wrong defendant, or wrong venue can derail a meritorious case.
		\end{itemize}
		\begin{alertblock}{Stop and reconsider hiring counsel if\ldots}
			\small Any criminal record; any allegation of fraud or misrepresentation; the case was \emph{denied} (different cause of action); class action or multiple plaintiffs; you are outside the U.S.; or you need in forma pauperis status.
		\end{alertblock}
	\end{frame}
	
	% ---------- Section 1 ----------
	\begin{frame}{Five Threshold Requirements}
		\small
		\begin{enumerate}
			\item \textbf{Clear, non-discretionary duty.} Adjudicating a properly filed application is non-discretionary; the \emph{outcome} is discretionary.
			\item \textbf{Clear right to performance.} You must be the actual applicant or petitioner.
			\item \textbf{No other adequate remedy.} Case status checks, service requests, and ombudsman inquiries are informal --- but try them first.
			\item \textbf{Unreasonable delay.} Measured by the six-factor \emph{TRAC} test; practically, well past published USCIS processing times.
			\item \textbf{Jurisdiction and venue.} 28 U.S.C. \S\S~1361, 1331; venue under \S~1391(e).
		\end{enumerate}
	\end{frame}
	
	\begin{frame}{The TRAC Six Factors}
		\small
		\emph{Telecommunications Research \& Action Center v. FCC}, 750 F.2d 70 (D.C. Cir. 1984)
		\begin{enumerate}
			\item Time governed by a ``rule of reason'' --- cite the agency's own published processing times.
			\item Congressional timetable --- 8 U.S.C. \S~1571(b) (180 days); \S~1447(b) for N-400s (120 days post-interview).
			\item Delays less tolerable when health and welfare are at stake --- plead concrete harms.
			\item Effect of expediting on competing agency priorities --- anticipate the ``line-jumping'' argument.
			\item Nature and extent of interests prejudiced --- family separation, work authorization, retrogression.
			\item No need to show impropriety or bad faith --- institutional sluggishness suffices.
		\end{enumerate}
	\end{frame}
	
	% ---------- Section 2 ----------
	\begin{frame}{Exhaust Administrative Remedies First}
		\small
		\begin{table}
			\centering
			\begin{tabular}{@{}ll@{}}
				\toprule
				\textbf{Step} & \textbf{What to Do} \\
				\midrule
				Case status & Screenshot the online case status; save as PDF. \\
				Service request & Contact Center (1-800-375-5283); 30-day response; tracking number. \\
				Ombudsman & Form DHS-7001 via the Ombudsman website; wait 60 days first. \\
				Congressional & Senator/Representative constituent services; keep all correspondence. \\
				Tier 2 & InfoPass / Tier 2 escalation; document date, officer, response. \\
				\bottomrule
			\end{tabular}
		\end{table}
		\begin{alertblock}{Build your exhibit packet}
			\small Gather every receipt, response, and screenshot chronologically; attach the most important as exhibits.
		\end{alertblock}
	\end{frame}
	
	% ---------- Section 3 ----------
	\begin{frame}{Choosing Your District Court}
		\small
		\textbf{Venue under 28 U.S.C. \S~1391(e) --- proper where:}
		\begin{enumerate}
			\item A defendant resides (federal officers often ``reside'' in D.C., or where the office is);
			\item A substantial part of the events occurred (service center / field office); or
			\item The plaintiff resides, if no real property is involved.
		\end{enumerate}
		\begin{block}{Practical choice}
			\small File where you live --- no travel, your local rules, easy reading.
		\end{block}
		\textbf{Gather from the court website:} Local Rules, Form JS-44, Form AO 440, CM-ECF pro se rules, pro se handbook.
	\end{frame}
	
	% ---------- Section 4 ----------
	\begin{frame}{Identifying the Correct Defendants}
		\small
		\begin{alertblock}{Common error}
			USCIS, DHS, and the State Department are \emph{not} proper defendants --- sue the \textbf{officers in their official capacity}.
		\end{alertblock}
		\textbf{For a USCIS delay (I-485, I-130, I-140, I-765, N-400, I-751):}
		\begin{itemize}
			\item Director of USCIS
			\item Secretary of Homeland Security
			\item Director of the Field Office / Service Center holding your file
			\item Attorney General of the United States
		\end{itemize}
		\textbf{Consular / State Department delay:} add the Secretary of State and the Consul/Consul General.
	\end{frame}
	
	% ---------- Section 5 ----------
	\begin{frame}{Anatomy of a Mandamus Complaint}
		\small
		FRCP 8: (1) grounds for jurisdiction, (2) short and plain statement of the claim, (3) demand for relief.
		\begin{enumerate}
			\item \textbf{Caption} --- court, parties, case no. blank, title.
			\item \textbf{Introduction} --- plain-language narrative (2--4 paragraphs).
			\item \textbf{Jurisdiction \& Venue} --- \S\S~1331, 1361, 5 U.S.C. \S~702; \S~1391(e)(1); exhaustion.
			\item \textbf{Parties} --- identify each, residence, receipt number.
			\item \textbf{Statement of Facts} --- chronological, numbered, exhibit references.
			\item \textbf{Causes of Action} --- Count I mandamus (\S~1361); Count II APA (\S\S~555(b), 706(1)).
			\item \textbf{Prayer for Relief} --- compel adjudication \emph{forthwith}, retain jurisdiction, costs/EAJA.
		\end{enumerate}
	\end{frame}
	
	% ---------- Section 6 ----------
	\begin{frame}{Supporting Documents}
		\small
		\begin{columns}[T]
			\begin{column}{0.5\textwidth}
				\textbf{Required Forms}
				\begin{itemize}
					\item \textbf{JS-44} Civil Cover Sheet
					\item \textbf{AO 440} Summons (one per defendant)
					\item \textbf{Proposed Order}
				\end{itemize}
			\end{column}
			\begin{column}{0.5\textwidth}
				\textbf{Typical Exhibits}
				\begin{itemize}
					\item A: I-797C Receipt Notice
					\item B: Underlying approval (I-130/I-140)
					\item C: Biometrics confirmation
					\item D: Service requests \& responses
				\end{itemize}
			\end{column}
		\end{columns}
		\vspace{1em}
		\begin{alertblock}{FRCP 5.2 redaction}
			\small SSN and account numbers: last four digits only. DOB: year only. Minors: initials only.
		\end{alertblock}
	\end{frame}
	
	% ---------- Section 7 ----------
	\begin{frame}{Filing Through PACER / CM-ECF}
		\small
		\begin{block}{PACER vs. CM-ECF}
			\textbf{PACER} = viewing court records. \textbf{CM-ECF} = electronic filing system.
		\end{block}
		\textbf{Paper filing package:}
		\begin{itemize}
			\item Signed original complaint + copies per defendant;
			\item Civil Cover Sheet (JS-44);
			\item Completed summonses (AO 440);
			\item Filing fee or IFP application.
		\end{itemize}
	\end{frame}
	
	% ---------- Section 8 ----------
	\begin{frame}{Filing Fees and In Forma Pauperis}
		\small
		\begin{block}{Current Fee}
			\$405 total (\$350 statutory filing fee + \$55 administrative fee).
		\end{block}
		\textbf{In Forma Pauperis (28 U.S.C. \S~1915):}
		\begin{itemize}
			\item File Form AO 240 if unable to afford fees.
			\item IFP triggers automatic court screening under \S~1915(e)(2).
		\end{itemize}
	\end{frame}
	
	% ---------- Section 9 ----------
	\begin{frame}{Serving the United States --- Rule 4(i)}
		\small
		\begin{block}{Rule 4(i) Service Requirements}
			Must send Summons and Complaint by certified mail to:
			\begin{enumerate}
				\item U.S. Attorney for your district;
				\item Attorney General (Washington, D.C.); and
				\item Each named official defendant.
			\end{enumerate}
		\end{block}
		\begin{itemize}
			\item \textbf{Deadline:} 90 days under Rule 4(m).
			\item \textbf{Proof of Service:} Always file green cards/receipts with the court.
		\end{itemize}
	\end{frame}
	
	% ---------- Section 10 ----------
	\begin{frame}{After You File: The 60-Day Window}
		\small
		FRCP 12(a)(2): The United States has \textbf{60 days from service} to respond.
		\begin{itemize}
			\item \textbf{Days 14--45:} A USAO attorney is assigned and often reaches out to discuss resolution.
			\item \textbf{Days 30--60:} Agency review occurs; most clean cases get adjudicated here.
			\item \textbf{If resolved:} File a Stipulation of Dismissal under FRCP 41(a)(1)(A)(ii).
		\end{itemize}
	\end{frame}
	
	% ---------- Section 11 ----------
	\begin{frame}{Ten Common Pro Se Mistakes}
		\small
		\begin{columns}[T]
			\begin{column}{0.5\textwidth}
				\begin{enumerate}
					\item Filing too soon
					\item Naming an agency as defendant
					\item Skipping Rule 4(i) service steps
					\item Failing to file proof of service
					\item Filing in the wrong venue
				\end{enumerate}
			\end{column}
			\begin{column}{0.5\textwidth}
				\begin{enumerate}\setcounter{enumi}{5}
					\item Seeking approval, not adjudication
					\item Forgetting FRCP 5.2 redaction
					\item Missing settlement outreach
					\item Ignoring local rules
					\item Missing deadlines
				\end{enumerate}
			\end{column}
		\end{columns}
	\end{frame}
	
	% ---------- Closing ----------
	\begin{frame}{Closing Note}
		\begin{block}{General Information Only}
			\small This presentation is for informational purposes only and does not constitute legal advice or create an attorney-client relationship.
		\end{block}
	\end{frame}
	
\end{document}